% ------------------------------------------------------------
% English working translation layer.
% Source file: letters/Letter_Manin_1987.tex
% ------------------------------------------------------------

\documentclass[11pt,a4paper]{article}
\usepackage[utf8]{inputenc}
\usepackage[T1]{fontenc}
\usepackage[english]{babel}
\usepackage{amsmath,amssymb,amsthm}
\usepackage{amsfonts}
\usepackage{mathtools}
\usepackage{geometry}
\usepackage{hyperref}
\usepackage{enumitem}
\usepackage{booktabs}
\usepackage{array}
\usepackage{mathrsfs}
\usepackage{xcolor}
\usepackage{microtype}

\geometry{a4paper, top=2.5cm, bottom=2.5cm, left=2.5cm, right=2.5cm}

\theoremstyle{plain}
\newtheorem{theorem}{Theorem}
\newtheorem{lemma}[theorem]{Lemma}
\newtheorem{proposition}[theorem]{Proposition}
\newtheorem{corollary}[theorem]{Corollary}
\newtheorem{conjecture}[theorem]{Conjecture}

\theoremstyle{definition}
\newtheorem{definition}[theorem]{Definition}
\newtheorem{example}[theorem]{Example}
\newtheorem{remark}[theorem]{Remark}

\DeclareMathOperator{\Ext}{Ext}
\DeclareMathOperator{\Hom}{Hom}
\DeclareMathOperator{\Ker}{Ker}
\DeclareMathOperator{\Coker}{Coker}
\DeclareMathOperator{\Spec}{Spec}
\DeclareMathOperator{\Proj}{Proj}
\DeclareMathOperator{\Gal}{Gal}
\DeclareMathOperator{\Aut}{Aut}
\DeclareMathOperator{\End}{End}
\DeclareMathOperator{\Ber}{Ber}
\DeclareMathOperator{\Res}{Res}
\DeclareMathOperator{\Id}{Id}
\DeclareMathOperator{\Sp}{Sp}

\begin{document}
\begin{center}\small\emph{Literal English translation working draft. Formulae, numbering, and TeX structure are preserved; this is not yet a modern recast.}\end{center}
\vspace{0.5em}


\begin{flushright}
Princeton, the 25 septembre 1987
\end{flushright}

\vspace{0.5cm}

\noindent Dear Manin,

\vspace{0.3cm}

Inspired by a conversation with Witten, I can now compactify --- with a divisor with normal crossing at infinity, the moduli stack of super Riemann surfaces. As I have to think, I continue in French.

\vspace{0.5cm}

\section*{1. Super-curves generiquement smooth, of Cohen-Macaulay}
\addcontentsline{toc}{section}{1. Super-curves generiquement smooth, of Cohen-Macaulay}

In this section, je propose a analogue en super-geometrie algebraic of the notion of curve geometriquement reduite.

The notion recue of super curve smooth $X/S$ is: superschema $X$ on the superschema $S$, smooth of dimension relative $(1,1)$, equipped with $D \subset T_{X/S}$ locally direct factor of type dimension $(0,1)$, and ``also not integrable that possible''. For a tel systeme $(X,D)$, one identifie $T_{X/S}/D$ a $D^{\otimes 2}$ by $D_1 \otimes D_2 \mapsto \frac{1}{2}[D_1,D_2]$. One can dualiser $D \hookrightarrow T_{X/S}$ en
\[
\Omega^1_{X/S} \twoheadrightarrow D^{\vee},
\]
and the isomorphism $\Omega^1_{X/S} \twoheadrightarrow T_{X/S}/D \sim D^{\otimes 2}$ se dualise en
\[
0 \longrightarrow D^{\vee \otimes 2} \longrightarrow \Omega^1_{X/S} \longrightarrow D^{\vee} \longrightarrow 0
\]
of or a isomorphism $\Ber(\Omega^1_{X/S}) \simeq D^{\vee}$. Of the:
\begin{equation}\label{eq1.2}
\delta : \Omega^1_{X/S} \longrightarrow \omega_{X/S} := \Ber(\Omega^1_{X/S}), \quad \text{i.e.} \quad \delta : \text{derivation} \quad \mathcal{O}_X \longrightarrow \omega_{X/S}.
\end{equation}

Mon choix of signes is can-be ambigu. For $X$ of coordonnees local $(z,\theta)$, $D = (\partial_\theta + \theta\partial_z)$, je veux
\[
\delta : f \longmapsto (\partial_\theta + \theta\partial_z)f \cdot [dz/d\theta]
\]
or $[dz/d\theta]$ is the base of $\Ber(\Omega^1_{X/S})$ defined by the base $(dz,d\theta)$ of $\Omega^1_{X/S}$.

Reciproquement, $\delta$ determine $D$: $D$ is the orthogonal of $\Ker(\delta)$, or the multiples of $\delta/\underline{b}$ for $\underline{b}$ the base of $\Ber(\Omega^1_{X/S})$. Noter that for $\delta : \Omega^1_{X/S} \to \omega_X$, epimorphique, it ne suffit not that $\Ker(\delta)^{\perp}$ let also not integrable that possible for $\delta$ provenir of a structure of supercourbe: the normalisation must be correcte: $\delta = D \cdot \underline{b}$, $\underline{b} = [D^2/\partial_z]^{-1}$.

En coordonnees local $(z,\theta)$, for $D = \partial_\theta + \theta\partial_z$ (a derivation impaire), one has

\begin{lemma}\label{lem1.1}
For that $\delta = (aD + bD^2)[dz/d\theta]$ definisse a structure of supercourbe, it is necessary and sufficient that $a$ let inversible and that $[(aD+bD^2)(aD+bD^2)] = [D^2/\partial_z]$, i.e.
\begin{equation}\label{eq1.1.1}
\Ber\begin{pmatrix} a^2 + b\partial_z^2 b + a\partial_z a & b \\ -(a\partial_z a + b\partial_z^2 a) & a \end{pmatrix} = 1, \quad \text{ou} \quad \Ber\begin{pmatrix} a & b \\ c & d \end{pmatrix} = (ad - \beta\gamma)d^{-2}.
\end{equation}
\end{lemma}

\noindent \textit{Proof}: $(aD+bD^2)^2 = (a\partial_z a + b\partial_z^2 a)D + (a^2 + b\partial_z^2 b + a\partial_z a)\partial_z$.

The cas which we servira the more is celui or $\delta = D \cdot \underline{b}$ is a structure of supercourbe, or $S_0 \hookrightarrow S$ is defined by a ideal of carre zero, and or $D_1 = D + a_1 D + b_1 D^2$ coincide with $D$ to the-dessus of $S_0$. The probleme se ramene then a supertrace:

\newpage

\begin{corollary}\label{cor1.2}
Under the hypotheses ci-dessus, the conditions of \ref{lem1.1} equivalent a:
\begin{equation}\label{eq1.2.1}
a_1 + Db_1 = 0.
\end{equation}
\end{corollary}

\noindent \textit{Proof}: ecrire $(2a_1 + Db_1) - a_1 = 0$.

Bien on, ci-dessus, $a$ and $a_1$ are pairs, $b$ and $b_1$ impairs.

\begin{remark}\label{rem1.3}
(i) For $X/S$ smooth of dimension relative $(1,1)$, let $\underline{b}$ a base of $\omega_X$ and $\delta = D \cdot \underline{b} : \mathcal{O} \to \omega$ a derivation. The condition that $D$ generated a direct factor local of dimension $(1,0)$ of $T_{X/S}$ equivaut a the surjectivite of $\delta : \Omega^1_{X/S} \to \omega_{X/S}$. For $S$ a scheme ordinaire (pair), it equivaut again a celle of the morphism induces by $\delta : \mathcal{O}_X^- \to \omega_X^-$.

(ii) For $\delta = D \cdot \underline{b} : \Omega^1_{X/S} \to \omega_{X/S}$ surjective, si, on a open schematiquement dense $\delta$ is a structure of supercourbe, then $\delta$ en is a partout: en terme of a structure of supercourbe of reference, the Berezinien (\ref{eq1.1.1}) a denominateur inversible and si $\Ber = 1$ on a open schematiquement dense, this identity vaut partout.
\end{remark}

\begin{definition}\label{def1.4}
A supercourbe on $S$ is a superschema $X/S$, plat, of Cohen-Macaulay, which on a open dense fiber a fiber is smooth of dimension relative $(1,1)$, equipped with $\delta$ of a derivation $\delta : \mathcal{O}_X \to \omega_{X/S}$ (or $\omega_{X/S}$ is the sheaf dualisant) such that

\noindent (i) on each fiber geometric, $\delta$ induces a isomorphism $\mathcal{O}_{X_s}^- \sim \omega_{X_s}^-$

\noindent (ii) on the open of smooth, $\delta$ is a structure of supercourbe smooth.
\end{definition}

of after \ref{rem1.3} (ii), the condition \ref{def1.4} (ii) can be remplacee by:

\noindent (\ref{def1.4}(ii)') On a open smooth on $S$ and schematiquement dense, $\delta$ defines a structure of supercourbe smooth on $S$.

\medskip

\noindent \textbf{1.5} For $S$ a scheme ordinaire, explicitons this definition en terme of the scheme $C := (X, \mathcal{O}_X^+)$ on $S$. The conditions plat, Cohen-Macaulay, generiquement smooth fiber a fiber, donnent

\noindent (1.5.1) $C/S$ is a curve plate a fibers geometriquement reduites.

The sheaf coherent $\mathcal{L} := \mathcal{O}_X^-$ is plat on $S$, relativement of Cohen-Macaulay, locally free of rank $1$ on a open dense fiber a fibers, and $\mathcal{O}_X^- \cdot \mathcal{O}_X^- = 0$.

For a tel superschema $X = (C, \mathcal{O}_C \oplus \mathcal{L})$, oux the sheaf dualisant $\omega_X$ se calcule as suit. One dispose of
\[
p : X \longrightarrow C
\]
and $\omega_{X/S} = p^!\omega_{C/S} = \underline{\Hom}_{\mathcal{O}_C}(\mathcal{O}_X, \omega_C)$.

the hypothese ``Cohen-Macaulay relative'' assure the nullite of the $\Ext$ superieurs, and that $\omega_{X/S}$ is plat on $S$, of formation compatible a all changement of base. one has
\[
@@MATH6@@
\]

one has $\mathcal{O}_X^- \cdot \omega_{X/S}^+ = 0$ and $\mathcal{O}_X^- \otimes \omega_{X/S}^- \to \omega_{X/S}^+$ is the evaluation $\mathcal{L} \otimes \Hom(\mathcal{L},\omega_C) \to \omega_C$.

\newpage

In the cas smooth, en coordonnees local $(z,\theta)$, $\omega_X$ admet the base $[dz/d\theta]$ and $\omega_X^-$ s'identifie a $\omega_C$ by
\[
\theta[dz/d\theta] \longmapsto dz.
\]

A derivation $\delta : \mathcal{O}_X \to \omega_X$ is defined by
\[
\begin{aligned}
\delta^+ &: \mathcal{O}_C \longrightarrow \omega_C && \text{(une derivation)} \\
\delta^- &: \mathcal{L} \longrightarrow \Hom(\mathcal{L}, \omega_C) && (\mathcal{O}_C\text{-lineaire})
\end{aligned}
\]

The conditions ( ) donnent $\delta xy = \cdots$ for $x$ or $y$ pair. For $x$ and $y$ impairs, it faudrait exprimer the symetrie of the form bilineaire $\mathcal{L} \otimes \mathcal{L} \to \omega_C$ defined by $\delta^-$. As $\mathcal{L}$ is generiquement of rank 1, c'is automatique.

the axiome \ref{def1.4}(i) given: $\delta^-$ is a isomorphism.

the axiome \ref{def1.4}(ii) signifie that, on the open of smooth,
\[
\delta^+ : \mathcal{O}_C \to \omega_C \quad \text{est} \quad d : \mathcal{O}_C \to \Omega^1_C = \omega_C. \quad \text{En effet}
\]
(a) this condition is verifiee for a supercourbe smooth en coordonnees local standard: one has $\delta z = \theta[dz/d\theta] \sim dz$.

(b) deux $\delta$ verifiant the condition are locally isomorphes for the topologie étale.

the argument (b) force en car. 2. Voici a more direct, which marche. Let en coordonnees local
\[
\delta f = (u(z)\partial_\theta + \theta v(z)\partial_z)f \cdot [dz/d\theta].
\]
one has $(u(z)\partial_\theta + \theta v(z)\partial_z)^2 = u(z)v(z)\partial_z$, and the condition
\[
[u(z)v(z)\partial_z / u(z)\partial_\theta + \theta v(z)\partial_z] = [\partial_z / \partial_\theta]
\]
given $v = 1$, as voulu.

\newpage

Parce that, on a open schematiquement dense, one has $\delta^+ = d$, c'is partout vrai. The datum on a scheme (by ex. !) of a supercourbe equivaut therefore a celle of $C$, $\mathcal{L}$ as en (1.5.1), more the datum of

\noindent (1.5.2) a form bilineaire (automatiquement symmetric) $\mathcal{L} \otimes \mathcal{L} \to \omega_C$ induisant a isomorphism $\mathcal{L} \xrightarrow{\sim} \Hom(\mathcal{L}, \omega_C)$.

As explique, one has

\noindent (1.5.3) $X = (C, \mathcal{O}_C \oplus \mathcal{L})$, $\omega_X = \omega_C \oplus \mathcal{L}$, with the structure of module defined by (1.5.2) and $\delta = (d, \Id)$.

En caracteristique 2, a family of supercourbes smooth n'is not necessarily locally (en haut) trivial. For this raison, je assumed desormais 2 inversible.

For that presqu'a present all marche on $\mathbb{Z}$, j'ai pris the definition of ``super-ring commutative'' or $a^2 = 0$ for $a$ impair, and j'ai fait gaffe a utiliser $D^2$ plutot that $[D,D]$ for the derivations impaires.

\medskip

\noindent \textbf{1.6} The description suivante of $\omega_{X/S}$ we will be utile.

Let $j : U \hookrightarrow X$ a open dense fiber a fiber and smooth. For every section of $j_*\omega_{U/S}$ and for $F$ a composante connected of $X-U$, finite on $S$, a residu $\Res_F(s) \in \mathcal{O}_S$ is defined. The sheaf $\omega_{X/S}$ is the under-sheaf of $j_*\omega_{U/S}$ form of the sections $s$ such that, locally for the topologie étale on $S$ and for $F$ as ci-dessus, all $f$ in $\mathcal{O}_X$, one ait $\Res_F(f \cdot s) = 0$.

On $S$ nilpotent, $\Res_F$ is a sum of residus on the branches passant by $F$ and, for each branche $B \hookrightarrow X$ (closed in $X$), is invariant by of the morphisms birationnels $B_1 \to B$. This permet of calculer $\omega_X$ a the Serre.

\begin{proposition}\label{prop1.7}
Let $X_0$ a supercourbe over a field $k$. A deformation infinitesimale of $X_0$ which is trivial en tant that deformation of superschema is trivial.
\end{proposition}

\noindent \textit{Proof}: It is enough to traiter the cas of deformations to the prime ordre (argument a the Schlessinger), and this cas se ramene to the cas of a deformation on $k[\epsilon]/(\epsilon^2)$ or on $k[\epsilon]$.

\medskip

\noindent a. \textbf{Deformation paire} (on $k[t]/(t^2)$). One passe to the point of vue 1.5. Let $X_0$ defined by $C$, $\mathcal{L}$ and $\mathcal{L} \xrightarrow{\sim} \Hom(\mathcal{L},\omega_C)$. A deformation a superschema constant can be vue as a deformation of $\mathcal{L} \xrightarrow{\sim} \Hom(\mathcal{L},\omega_C)$. A telle deformation is of the type $\varphi \circ a$, $a$ automorphism infinitesimal of $\mathcal{L}$. A tel automorphism infinitesimal $a = 1 + a_1$, admet the racine carree $1 + \frac{1}{2}a_1$, and $\varphi$ is conjugue of $\varphi_0$ by $(1+a_1)^{-1/2}$: it is enough to the verify on the lieu smooth, or $a$ is the multiplication by a element of $\mathcal{O}_{C \otimes k[t]/(t^2)}^*$ egal a 1 modulo $t$. Conjuguer $\varphi$ ne change not the class of isomorphie of deformation of supercourbe, and the deformation is therefore trivial.

\newpage

\noindent b. \textbf{Deformation impaire} (on $k[\epsilon]$)

A deformation of $(X_0, \delta_0)$, trivial as deformation of superschema, is of the type $(X,\delta)$ with $X$ deduces of $X_0$ by extension of scalars and of derivation structurale
\[
\delta = \delta_0 + \epsilon D \qquad (\delta_0 : \mathcal{O}_X \to \omega_X \text{ deduit de } \delta_0 \text{ initial par extension des scalaires}),
\]
or $D$ is a derivation impaire $\mathcal{O}_{X_0} \to \omega_{X_0}$. It faut prouver that si $\delta$ satisfies \ref{def1.4}(ii), then $\delta$ is conjugue of $\delta_0$ by a automorphism of $X$ trivial on $X_0$. A tel automorphism is of the form $f \mapsto f + \epsilon Ef$, for $E$ a derivation impaire of $\mathcal{O}_{X_0}$ in $\mathcal{O}_{X_0}$.

Of the point of vue 1.5, a derivation impaire $D : \mathcal{O}_{X_0} \to \omega_{X_0}$ is given by ses composantes
\[
@@MATH7@@
\]

For $\delta_0 + \epsilon D$ verifiant \ref{def1.4}(ii), $D_1$ is determine by $D_0$: it is enough to the verify on the lieu smooth, and the one has en coord. local
\[
Df = (u(z)\partial_z + v(z)\theta\partial_\theta)f \cdot [dz/d\theta],
\]
and $D_0$ is $f \mapsto u(z)\partial_z f \cdot [dz/d\theta] : \mathcal{O}_{X_0}^+ \to \omega_{X_0}^-$. Ecrivant that $D = \theta v(\partial_\theta + \theta\partial_z) + u\partial_z$ and appliquant \ref{cor1.2}, one sees that \ref{def1.4}(ii) signifie: $v = \partial_z u$:

$D_0$ determine $D$.

\newpage

A derivation impaire $E$ of $\mathcal{O}_X$ se decrit of the point of vue 1.5 by ses composantes
\[
@@MATH8@@
\]

One va prendre for $E$ the derivation of composantes $(D_0, 0)$. To verify that the automorphism infinitesimal $\Id + \epsilon E$ defined by $E$ conjugue $\delta_0$ en $\delta$, it is enough to the verify on the lieu smooth, and it is enough to verify that the composante ``0'' of $\mathcal{L}_E(\delta_0)$ is $D_0$. En coordonnees local, one has
\[
\begin{aligned}
E &: u(z)\theta\partial_z \\
\delta_0 &: (\partial_\theta + \theta\partial_z) \cdot [dz/d\theta] \\
\mathcal{L}_E(\delta_0) &= [E, \partial_\theta + \theta\partial_z] \cdot [dz/d\theta] + (\partial_\theta + \theta\partial_z) \mathcal{L}_E([dz/d\theta]) \\
[E, \partial_\theta + \theta\partial_z] &= [\partial_\theta + \theta\partial_z, u(z)\theta\partial_z] = u(z)\partial_z
\end{aligned}
\]
and the second terme contribue a the composante ``1'' of $\mathcal{L}_E(\delta_0)$.

one has retrouve the $D_0$ voulu.

\begin{remark}\label{rem1.8}
J'ai enfin compris this that signifie the condition (ii) of the definition \ref{def1.4}: that $\delta$ let not proper transpose (a signe a ajouter pres).
\end{remark}

\newpage

\section*{2. Curves stables}
\addcontentsline{toc}{section}{2. Curves stables}

\begin{definition}\label{def2.1}
A super curve stable $X/S$ is a super curve proper on $S$ (\ref{def1.4}) of which the fibers geometriques $X_s$ have of the reduites $(X_s, \mathcal{O}_{X_s}^+)$ which are of the curves stables.
\end{definition}

For $S$ spectre of a field algebriquement clos $k$, it is easy a the aide of 1.5 of determiner (locally for the topologie étale) the singularities possibles of the super curves stables.

On the singularity $xy = 0$, there is indeed exactly deux types of isomorphie of sheaf coherent of Cohen-Macaulay generiquement of rank 1. For $C$ the curve, and $B_1$, $B_2$ ses deux branches, this are $\mathcal{O}_C$ and $\mathcal{O}_{B_1} \oplus \mathcal{O}_{B_2}$. The deux sheaves are isomorphes a leur dual (a value in $\mathcal{O}_C$, or $\omega_C$, that revient to the same). Moreover, all automorphism of a of these sheaves admet (locally étale) a racine carree, and 1.5 given deux types possibles of singularities of super-curves. Voici a description by coordonnees and equations of the deux possibilites. One given: coordonnees / equations / description of $\omega_X$ as under-sheaf of the direct image of sa restriction to the lieu smooth / description of $\delta$.

\medskip

\noindent \textbf{2.2. I.} coordonnees $z_1, z_2, \theta$

equation $z_1 z_2 = 0$

complement of the point singular: deux branches, with the coordonnees
\[
@@MATH9@@
\]

$\omega_X$: free of base $\underline{b}$, datum on each branche by
\[
\underline{b} = @@MATH10@@
\]

$\delta$: on each branche: $f \mapsto @@MATH11@@$

One denotes $D$ the derivation datum on each branche by: $@@MATH12@@$ $D : \mathcal{O} \to \mathcal{O}$, impair.

\medskip

\noindent \textbf{II.} coordonnees $z_1, z_2, \theta_1, \theta_2$

equations $z_1 z_2 = z_1\theta_2 = \theta_1 z_2 = \theta_1\theta_2 = 0$

complement of the point singular: deux branches, coordonnees $z_1, \theta_1$ ($z_2 = \theta_2 = 0$) and $z_2, \theta_2$ ($z_1 = \theta_1 = 0$)

$\omega_X$: generated by the forms data as suit on each branche:
\[
\begin{aligned}
&[dz_1/d\theta_1] := @@MATH13@@ \qquad [dz_2/d\theta_2] := @@MATH14@@ \\
&\text{et } \frac{\theta_1}{z_1}[dz_1/d\theta_1] - \frac{\theta_2}{z_2}[dz_2/d\theta_2] := @@MATH15@@
\end{aligned}
\]

$\delta$: on each branche: $(\partial_{\theta_i} + \theta_i\partial_{z_i})f \cdot [dz_i/d\theta_i]$

One sets $\underline{b} := [dz_1/d\theta_1] + [dz_2/d\theta_2]$ and one denotes $D$ the derivation impaire datum en dehors of the point singular by

\newpage

$D := \partial_{\theta_i} + \theta_i\partial_{z_i}$ on each branche. C'is a derivation rational: it n'envoie not $\mathcal{O}$ in $\mathcal{O}$: one has
\[
\theta_i \longmapsto @@MATH16@@
\]
but $\delta = D \cdot \underline{b}$ envoie $\mathcal{O}$ in $\omega_X$.

In I and II, the calcul of $\omega_X$ can se faire as en 1.6. Let $s_i$ of the sections of $\omega_{X/S}$ on $(B_i - \text{point singular})$ ($i = 1,2$); for that the $s_i$ proviennent of a section of $\omega_X$, it is necessary and sufficient that for every $f$ in $\mathcal{O}_X$, induisant $f_i$ on the branche $B_i$, one ait $\Res_0(f_1 s_1) + \Res_0(f_2 s_2) = 0$. Je recalls that $\Res_0(\theta/z \cdot [dz/d\theta]) = 1$.

In the cas I, each branche (closed) $B_i$ can be vue as a curve equipped with a structure of spin ramifiee en 0: a racine carree of $\omega_{C_i}(0)$ [$\zeta_i = B_i$ red.; pole simple permis en 0]. One recolle $B_1$ and $B_2$ selon the feuille integrable $z_i = 0$ of the distributions $\partial_\theta \pm \theta(z_i\partial_{z_i})$.

In the cas II, each branche (closed) $B_i$ can be vue as a supercourbe smooth; they are recollees by a point.

\newpage

\noindent \textbf{2.3} We have a considerer the deformations suivantes of the super curves I and II. The deformations are on $k[t]$.

\medskip

\noindent \textbf{I.} coordonnees $z_1, z_2, \theta$

equation $z_1 z_2 = t$ \hfill \begin{minipage}{5cm}\small in the plan $(z_1,z_2)$, the derivation $z_1\partial_1 - z_2\partial_2$ annule $t$\end{minipage}

The complement of the point singular $(0,0)$ ($t=0$) is recouvert by the deux cartes:
\begin{itemize}
\item[(1)] $z_1$ inversible, coordonnees $(z_1, \theta)$, $z_2 = t/z_1$
\item[(2)] $z_2$ inversible, coordonnees $(z_2, \theta)$, $z_1 = t/z_2$
\end{itemize}

$\omega_{X/S}$ is free of base $\underline{b}$ given by
\[
\begin{aligned}
\underline{b} &= [dz_1 \wedge dz_2 / d\theta] / [d(z_1 z_2)] \\
&= @@MATH17@@
\end{aligned}
\]

$\delta$ is given in each carte by
\begin{itemize}
\item[(1)] $(\partial_\theta + \theta z_1\partial_{z_1}) \cdot \underline{b}$
\item[(2)] $(\partial_\theta - \theta z_2\partial_{z_2}) \cdot \underline{b}$
\end{itemize}

Puisque $\delta$ envoie $\mathcal{O}_X^-$ in $\omega_{X/S}$ en dehors of $(0,0)$, by profondeur, c'is egalement vrai en $(0,0)$. The platitude is clear: intersection complete relative.

\medskip

\noindent \textbf{II.} coordonnees $z_1, z_2, \theta_1, \theta_2$

equations $z_1 z_2 = -t^2$, $z_1\theta_2 = t\theta_1$, $z_2\theta_1 = -t\theta_2$, $\theta_1\theta_2 = 0$

cartes (reunion: $X - (0,0)$):
\begin{itemize}
\item[(1)] $z_1$ inversible, coord. $z_1, \theta_1$, $@@MATH18@@$
\item[(2)] $z_2$ inversible, coord. $z_2, \theta_2$, $@@MATH19@@$
\end{itemize}

\newpage

Si $A$ is the ring of coordonnees of the space total, one has $A^+ = k[z_1, z_2, \sqrt{z_1 z_2}]$; one can see $A^+$ as the under-ring of $k[\sqrt{z_1}, \sqrt{z_2}]$ of the polynomes of degree total integer. By $\theta_1 \mapsto \sqrt{z_1}$, $\theta_2 \mapsto \sqrt{z_2}$, one can then identifier $A^-$ to the $A^+$-module of the sums of monomes of degree total $\frac{1}{2}$-integer not integer. This rend clear that $t = \sqrt{z_1 z_2}$ n'is not diviseur of 0, this the platitude of $X$ on $k[t]$.

Voici of the sections of $\omega_{X/S}$ en dehors of 0, ecrites in the cartes (1) and (2)
\[
@@MATH20@@ \qquad
@@MATH21@@ \qquad
@@MATH22@@
\]

A argument of profondeur shows that these sections se prolongent (uniquement) en of the sections on $X$. They relevent of the generateurs of the $\omega$ of the fiber $t=0$. Puisque $\omega_{X/S}$ is plat on $S$ of formation compatible to the changement of base, they engendrent $\omega_{X/S}$.

Enfin, $\delta$ is given in each carte by
\begin{itemize}
\item[(1)] $\delta f = (\partial_{\theta_1} + \theta_1\partial_{z_1})f \cdot [dz_1/d\theta_1]$
\end{itemize}
One verifies the compatibility in the intersection of the cartes.

\newpage

\begin{remark}\label{rem2.4}
The morphism $\delta : \Omega^1_{X/S} \to \omega_{X/S}$ is given by
\[
@@MATH3@@
\]; son image is generated by $[dz_1/d\theta_1]$ and $[dz_2/d\theta_2]$. Son kernel is of longueur 1, concentre en $(0,0)$. Indeed, si one $Y \cdot \theta_1/z_1 [dz_1/d\theta_1] = -\theta_2/z_2 [dz_2/d\theta_2]$ by $\theta_1, \theta_2, z_1$ or $z_2$ one trouve respectivement $0, 0, \theta_1[dz_1/d\theta_1]$ and $\theta_2[dz_2/d\theta_2]$.
\end{remark}

\begin{theorem}\label{thm2.5}
The deformation 2.3 is the deformation verselle (singular of) of the supercourbe 2.2.
\end{theorem}

It s'agit plutot of completer 2.2 and 2.3 a the origine and of deformation verselle formelle. Bien on, it s'agit of deformation as supercourbe, not only as superschema.

\begin{theorem}\label{thm2.6}
\begin{itemize}
\item[(i)] A supercourbe stable n'a not of automorphism infinitesimal not trivial.
\item[(ii)] The champ of the supercourbes stables of genre $g \geq 2$ is a champ algebraic proper and smooth.
\item[(iii)] It compactifie the champ of the supercourbes smooth of genre $g$ with a the infini a diviseur a croisements normaux.
\end{itemize}
\end{theorem}

A vrai dire, je n'ai satisfies 2.6 that modulo the superanalogue of the theorem of representabilite of Artin, and the superanalogue of the theorem of prorepresentabilite formelle of Schlessinger.

\newpage

Of the arguments standard ramement 2.6 (ii),(iii) a 2.6 (i) and 2.5 + verification of the critere valuatif of proprete. For this dernier, one can ne considerer that of the schemes of base and passe to the langage 1.5. Enfin, for prouver 2.5, it is enough to the faire to the prime ordre: si 2.3 on $k[t]/(t^2)$ is the seule deformation to the prime ordre, the space tangent a the space of the deformation verselle is of dimension $(1,0)$ and every deformation on $k[t]$ which prolonge celle on $k[t]/(t^2)$ is verselle.

C'is en examinant the critere valuatif of proprete that je suis tombe on the deformations 2.3. Of this point of vue, they ne are not surprenantes. The necessite of the type II vient of this that si a curve of genre $g$ degenere en deux curves of genre $g_1$ and $g_2$ attached en a point, then $\omega_{X_0}$ n'a not of racine carree, car of degree impair on each composante.

The diviseurs a the infini are caracterises by the fait that, on eux, a of the points singuliers admet (locally for the topologie étale) a equation 2.2. The remark 2.4 fournit to the-dessus of chacun of eux a section ``the point singular'' of image (locally) the support of $\Coker(\delta : \Omega^1_{X/S} \to \omega_{X/S})$. The diviseurs a the infini are of the under-varieties of codimension $(1,0)$.

J'avoue n'avoir of 2.5 that a proof degueulasse: the liasse of calculs which suit. C'is a peu a brouillon but, mettant to the net, je ne ferais guere mieux.

\newpage

\section*{3. Deformations to the prime ordre}
\addcontentsline{toc}{section}{3. Deformations to the prime ordre}

\noindent \textbf{3.1} For a moment, one va oublier $\delta$ and calculer the deformation verselle, to the $1^{\text{er}}$ ordre, of the singularity I or II. The calcul is penible, but standard. For a superschema a singularity isolee quelconque, it faut

(a) ecrire $X = \Sp(A)$, $A = k[\underline{x}, \underline{\theta}] / \text{equations } f_i$

(b) trouver a systeme of syzygies $S_k$: $\sum a_i(\underline{x},\underline{\theta}) f_i = 0$ the engendrant all modulo the syzygies triviales,
\[
f_i f_j - f_j f_i = 0 \quad \text{(cas pair ou $i,j$ impair)}
\]
\[
f_i f_j + f_j f_i = 0 \quad \text{($f_i$ et $f_j$ impairs)}
\]

(c) calculer the cohomology of the complex of $A$-modules libres
\[
\bigoplus \Bigl(A \cdot \frac{\partial}{\partial x_i} \text{ et } A \cdot \frac{\partial}{\partial \theta_j}\Bigr) \xrightarrow{\frac{\partial f}{\partial x}} \bigoplus A \cdot \frac{\partial}{\partial f_i} \xrightarrow{\text{coeffs de } S} \bigoplus A \cdot \frac{\partial}{\partial S_k}
\]
the $H^1$ parametre the deformations.

\medskip

\noindent \textit{Cas I}: $(z_1, z_2, \theta)$, $z_1 z_2 = 0$, not of syzygie.

Deformation universelle to the $1^{\text{er}}$ ordre:
\[
z_1 z_2 = t + \tau \theta \qquad (\text{sur } k[t,\tau]/(t^2, t\tau))
\]

The cas II is a other paire of manches.

\newpage

\noindent \textit{Cas II}: $z_1, z_2, \theta_1, \theta_2$ \hfill $(*)$
\[
z_1 z_2 = z_1 \theta_2 = \theta_1 z_2 = \theta_1 \theta_2 = 0 \qquad (**)
\]

syzygies:
\[
@@MATH23@@ \qquad (***)
\]

It s'agit of resoudre a systeme of 4 equations (a by syzygie) a 4 inconnues (a by equations $(**)$) in the ring $A$, exprimant that these inconnues satisfy the syzygies. It s'agit of calculer the space of the solutions modulo solutions triviales (obtained by derivation of the equations $(**)$).

Je note $\oplus\oplus$, $\oplus\ominus$, $\ominus\oplus$ and $\ominus\ominus$ the inconnues. The equations, corresponding to the syzygies, have for coefficients (solutions triviales:

\begin{center}
\begin{tabular}{c|cccc}
& $\oplus\oplus$ & $\oplus\ominus$ & $\ominus\oplus$ & $\ominus\ominus$ \\
\hline
$\theta_1$ & & $-z_1$ & & \\
$\theta_2$ & $-z_2$ & & & \\
& $-\theta_1$ & & $z_1$ & \\
& & $\theta_2$ & $z_2$ &
\end{tabular}
\qquad
\begin{tabular}{c|cccc}
& $(\partial_{z_1})$ & $(\partial_{z_2})$ & $(\partial_{\theta_1})$ & $(\partial_{\theta_2})$ \\
\hline
$\oplus\oplus$ & $z_2$ & $z_1$ & & \\
$\oplus\ominus$ & & $\theta_2$ & & $z_1$ \\
$\ominus\oplus$ & & & $\theta_1$ & $z_2$ \\
$\ominus\ominus$ & & & $\theta_2$ & $-\theta_1$
\end{tabular}
\end{center}

(each ligne, a syzygie) \hfill (each colonne, a relation trivial)

Let $(a,b,c,d)$ a solution. The modifiant by the deux dernieres solutions triviales, one can supposer $d$ pair. The equations $z_1 d = \pm \theta_1 b$, $z_2 d = -\theta_2 c$, donnent then $d = 0$. The combinaisons lineaires with $d = 0$ of solutions triviales are the combinaisons lineaires of $1^{\text{re}}$, $2^{\text{e}}$, $3^{\text{e}} \times \theta_2$, $4^{\text{e}} \times \theta_1$. This we ramene to the systeme

\newpage

\begin{center}
\begin{tabular}{c|ccc}
& $\oplus\oplus$ & $\oplus\ominus$ & $\ominus\oplus$ \\
\hline
$\theta_1$ & $-z_1$ & & \\
$\theta_2$ & $-z_2$ & & \\
& $-\theta_1$ & & $z_1$ \\
& & $\theta_2$ & $z_2$
\end{tabular}
\qquad ($\ominus\ominus = 0$)
\end{center}
\begin{center}
\begin{tabular}{c|cccc}
& $\oplus\oplus$ & $\oplus\ominus$ & $\ominus\oplus$ & $\ominus\ominus$ \\
\hline
$\partial_{z_1}$ & $z_2$ & & & \\
$\partial_{z_2}$ & $z_1$ & $\theta_2$ & & \\
$\partial_{\theta_1}$ & & & $\theta_1$ & $\theta_2$ \\
$\partial_{\theta_2}$ & & $z_1$ & $z_2$ & $-\theta_1$
\end{tabular}
\end{center}

\hfill Let solution $(a,b,c)$ ($d = 0$)

the equation $\theta_1 a = z_1 c$ impose that, modulo $(\theta_1, z_1)$, $(a,c)$ is multiple of the $2^{\text{e}}$ solution trivial. In particular, $a \in (z_1, \theta_1, z_2, \theta_2) = (z_1, \theta_1) \oplus (z_2, \theta_2)$. Similarly for $b$ and $c$: $a = a_1 + a_2$, $b = b_1 + b_2$, $c = c_1 + c_2$. Modulo the deux premieres solutions triviales, one obtains $a = 0$, $b_2 = 0$, $c_1 = 0$. the equation $\theta_1 b = z_1 a$ given then $b = \theta_1$,
\[
\begin{aligned}
a_1 \in (\theta_1), &\quad a_2 \in (\theta_2), \quad b_2 = c_1 = 0 \quad \text{et, par les dernieres equations} \\
& a_1 \in (\theta_1), \quad a_2 \in (\theta_2), \quad b_1 \in (\theta_1), \quad c_2 = 0 \\
& \qquad\qquad b_2 = 0 \qquad c_2 \in (\theta_2)
\end{aligned}
\]
and enfin, modulo the $\theta$-multiples of the premieres solutions triviales, every solution se ramene a the a of the suivantes:

\begin{center}
\begin{tabular}{|c|ccc|}
\hline
& $a$ & $b$ & $c$ \quad $d$ \\
\hline
$\theta_1$ & & & \\
$\theta_2$ & & & \\
& $\theta_1$ & & \\
& & $\theta_2$ & \\
\hline
\end{tabular}
\end{center}

and the deformation verselle to the $1^{\text{er}}$ ordre is

parametres $t_1, t_2, \tau_1, \tau_2$

coordonnees $z_1, z_2, \theta_1, \theta_2$

equations
\[
@@MATH24@@
\]

\newpage

\noindent \textbf{3.2} of after 3.1, the space versel of modules of deformations (singularity of) of a supercourbe is a under-space of the space versel of modules of the superschema under-jacent. To the prime ordre, one has

\begin{lemma}\label{lem3.3}
the space versel of modules of a singularity of super-curve is, to the $1^{\text{er}}$ ordre, the under-space suivant of the space versel of modules (calcule en 3.1) of the singularity of superschema:
\[
\begin{aligned}
\text{I} &: \tau = 0 \\
\text{II} &: \tau_1 = \tau_2 = 0, \; t_1 + t_2 = 0.
\end{aligned}
\]
\end{lemma}

It s'agit of savoir quand one can deformer $\delta$, en respectant the conditions suivantes:

(a) $\delta$ is with values in $\omega_X$

(b) on the lieu smooth, $\delta$ satisfies \ref{def1.4}(ii)

In the deformations considered (parametrisees by a under-super-scheme of the scheme which parametre the deformation 3.1), the complement $U$ of the point singular is reunion of deux branches:

$\odot$: $z_1$ inversible, $\odot$: $z_2$ inversible, admettant the coordonnees respectives $(z_1, \theta_1)$ and $(z_2, \theta_2)$. One decrit a section $w$ of $\omega_X$ by ses restrictions $w_1$ and $w_2$ to the deux branches. The condition on $w_1$ and $w_2$ is that for every section $f$ of $\mathcal{O}_X$, decrite by ses restrictions $f_1$ and $f_2$ to the deux branches, one ait in $\mathcal{O}_S$
\[
\Res_0(f_1 w_1) + \Res_0(f_2 w_2) = 0 \qquad (\text{cf. 1.6})
\]

Puisque $\omega_X$ is plat on $S$, for obtenir a set of generateurs of $\omega_X$, it is enough to relever a set of

\newpage

generateurs of $\omega_{X_0}$ ($X_0$ fiber speciale).

In the cas I (resp II), one denotes $\underline{b}$ the section suivante of $\omega_X$ on $U$:
\[
@@MATH4@@
\]

$D_0$ the derivation impaire on $U$
\[
\text{I} : @@MATH25@@ \qquad
\text{II} : @@MATH26@@
\]
and $\delta_0 = D_0 \cdot \underline{b}$. The $\delta_0$ n'envoie not necessarily $\mathcal{O}_X$ in $\omega_X$ and, for that a deformation as supercourbe let possible, it faut pouvoir corriger this $\delta_0$ for that it envoie $\mathcal{O}_X$ in $\omega_X$.

Distinguons the cas I and II, and the deformations paires and impaires.

\medskip

\noindent \textit{Cas I impair} \hfill $z_1 z_2 = \tau \theta$

Ici, $\underline{b}$ is in $\omega_X$ (one has $\underline{b} = [dz_1 dz_2/d\theta] / [d(z_1 z_2 - \tau\theta)]$), therefore is a base of $\omega_X$. It faudrait corriger the derivation impaire $D_0$ for that it let with values in $\mathcal{O}$. Decrivant a section of $\mathcal{O}_X$ by ses restrictions to the branches, one has
\[
z_1 = @@MATH27@@ \xrightarrow{D_0}
@@MATH28@@ \quad \text{qui n'est pas dans } \mathcal{O} :
\]
$@@MATH29@@$ $@@MATH30@@$ ne the is not, car $0$ n'is not in $\mathcal{O}_{X_0}$.

The derivation $D$ devrait envoyer $z_1$ on a quantite $\neq 0$ on the $2^{\text{e}}$ branche: impossible.

\newpage

\noindent \textit{Cas II. Absence of deformation impaire}

$X_0$ is bihomogene: one can attribuer the weight
\[
\theta_1 : (1,0), \; z_1 : (2,0) \qquad \theta_2 : (0,2), \; z_2 : (0,2)
\]
The deformation verselle of the singularity the is also:
\[
\text{poids} : \quad z_1 : (1,2) \quad z_2 : (2,1) \quad t_1 \text{ et } t_2 : (1,1)
\]
Of the lors, s'there exists a deformation with $\tau_1 = a\tau$, $\tau_2 = b\tau$, $(a,b) \neq (0,0)$, it en exists also a with $(a,b)$ remplace by $(\lambda a, \mu b)$. Si $a$ or $b = 0$, by symetrie, one has a deformation a $\tau_2 = 0$. Otherwise, one has deux deformations distinctes and $\tau_1$ and $\tau_2$ are en fait libres. It is enough to etudier the deformation on $k[\tau_i]$,
\[
z_1 z_2 = \tau_i \theta_1 \theta_2 \qquad z_1 \theta_2 = z_2 \theta_1 = \theta_1 \theta_2 = 0
\]

Calcul of $\omega_X$: one ecrit a element of $\omega_X$ of after this that it is on each branche:
\[
@@MATH31@@
\]
It s'agit of relever the generateurs:

$[dz_1/d\theta_1]$: not in $\omega_X$ car
\[
z_2 \cdot [dz_1/d\theta_1] = @@MATH32@@
\]
se releve en: $@@MATH33@@$

$[dz_2/d\theta_2]$: se releve en $@@MATH34@@$

$\frac{\theta_1}{z_1}[dz_1/d\theta_1] - \frac{\theta_2}{z_2}[dz_2/d\theta_2]$: se releve en idem on each branche.

\newpage

Cherchons a relever $\delta = D_0 \cdot \underline{b}$ (as en 3.1 II).

The lifting rational given by the same formule on each branche n'is not with values in $\omega_X$: one has
\[
\theta_i = @@MATH35@@ \longrightarrow @@MATH36@@
\]
and $D_0$ devrait be corrige by a derivation rational paire envoyant $\theta_i \to \frac{\theta_2}{z_2^2}[dz_2/d\theta_2]$ + replie on $2^{\text{e}}$ branche,

this which n'is not possible (respect of the support).

\medskip

\noindent \textit{Cas II, deformations paires}
\[
\begin{aligned}
z_1 z_2 &= 0 \qquad & z_1 \theta_2 = t_1 \theta_1 \qquad & \theta_1 \theta_2 = 0 \\
& & z_2 \theta_1 = t_2 \theta_2
\end{aligned}
\]

branche 1: $z_2 = 0$, $\theta_2 = t_1 \theta_1/z_1$ \hfill (coordonnees $z_1, \theta_1$)

branche 2: $z_1 = 0$, $\theta_1 = t_2 \theta_2/z_2$ \hfill (coordonnees $z_2, \theta_2$)

Cherchons again of the generateurs of $\omega_X$, under the form
\[
@@MATH37@@ \qquad u_i \in k[z_i, z_i^{-1}, \theta_i]
\]

Relever $[dz_1/d\theta_1]$: $\times \; \theta_2 = t_1 \theta_1/z_1 =$ residu $t_1$

releve en $[dz_1/d\theta_1] - t_1 \frac{1}{z_2} [dz_2/d\theta_2]$

$[dz_2/d\theta_2]$: se releve en $@@MATH38@@$

$@@MATH39@@$ se releve en the same.

\newpage

Let $D_0$ as precedemment. Matrice. one has
\[
@@MATH5@@
\]

Si $t_1 + t_2 \neq 0$, it faut corriger $D_0$ en $D_0 + \varepsilon D$, $D$ pair, and, a new, $D$ devrait envoyer $\theta_1$ on a quantite not zero on the $2^{\text{e}}$ branche: impossible.

One ne can therefore deformer that for $t_1 + t_2 = 0$.

That the deformations not exclues are possibles resulte of 2.3.

\newpage

\section*{4. Automorphismes infinitesimaux: proof of 2.6 (i)}
\addcontentsline{toc}{section}{4. Automorphismes infinitesimaux}

\begin{theorem}
A super-curve stable n'a not of automorphism infinitesimal.
\end{theorem}

Locally, on the lieu smooth, the sheaf of the automorphismes infinitesimaux
\[
\mathcal{A} \hookrightarrow T_X
\]
$*$ se projette isomorphiquement on $T_X / \text{distr. structure}$.

For $X : C, \mathcal{O}_C \oplus \omega^{1/2}$
\[
\mathcal{A} : \quad T_C \oplus \omega^{1/2} \otimes T_C.
\]

The prime travail is, pres of a singularity, of calculer quelles sections meromorphes of $\mathcal{A}$ are of the automorphismes infinitesimaux. It s'agit of envoyer $\mathcal{O}$ in $\mathcal{O}$.

\medskip

\noindent \textit{Cas I}: ici, $\mathcal{O}_X^-$ is a racine carree of $\omega_C$.

Let $D$ as en 3.1 I (one denotes $D_0$). A automorphism infinitesimal s'ecrit $aD + bD^2$ with $(-1)^{\deg a} \partial_z a + Db = 0$:
\[
a = \pm \frac{1}{2} \cdot (-1)^{\deg b} \cdot Db.
\]

For that it envoie $\mathcal{O}$ in $\mathcal{O}$, it faut:

\hspace{1cm} derivation paire: $b$ pair: $b \in \mathcal{O}^+$

\hspace{1cm} impaire: $b$ impair: $b \in \mathcal{O}^-$

\begin{center}
\fbox{$\mathcal{A} : \omega_C^{-1} \oplus \omega_C^{1/2} \oplus \omega_C^{-1}$}
\end{center}

\vspace{0.5cm}

$*$ car $[D, X + aD] \bmod \mathcal{O} = DX + (-1)^a a\partial_z$: unique a $a$ for amener in $\mathcal{O}$

\newpage

\noindent \textit{Cas II}: Similarly, one prend $D$ as en 3.1 II and one ecrit the derivations
\[
\begin{aligned}
\text{paire} &: \quad \tfrac{1}{2} Db \cdot D + bD^2 && (b \text{ pair}) \\
\text{impaire} &: \quad -\tfrac{1}{2} Db \cdot D + bD^2 && (b \text{ impair})
\end{aligned}
\]

Condition for that $b$, section rational of $\mathcal{O}$, given lieu a derivation: $\mathcal{O} \to \mathcal{O}$
\[
\begin{aligned}
\text{pair} &: \quad b \in \mathcal{O}^+, \text{ nul en } 0 \\
\text{impair} &: \quad b \in \mathcal{O}^-, \text{ nul en } 0
\end{aligned}
\]
\hfill en terme of the branches $B_i$

\[
\mathcal{A} : \mathcal{A}_1 \oplus \mathcal{A}_2, \qquad \mathcal{A}_i = T_{B_i}(\text{nul en } 0) \oplus T_{B_i}(\text{nul en } 0) \otimes \omega_{B_i}^{1/2}.
\]

A automorphism $\omega^{ab}$ pair / impair fournit therefore on the normalisee of a composante $A$ of $C$ a section of a sheaf en lines, of degree:
\[
\begin{aligned}
\text{pair} &: (2 - 2g) - \tfrac{1}{2}(\# \text{ points type I}) - (\# \text{ points type II}) \\
\text{impair} &: (1 - 2g) - \tfrac{1}{2}(\# \text{ points type I}) - (\# \text{ points type II})
\end{aligned}
\]

Under the hypothese ``stable'', these degrees are $< 0$.

\newpage

\section*{5. Critere valuatif of property}
\addcontentsline{toc}{section}{5. Critere valuatif of property}

Ici, one n'a as spaces of parametres a considerer that of the traits $(S,\eta,s)$, $S = \Spec(V)$ ($V$ field of the fractions $K$) and one can therefore utiliser the point of vue 1.5. of after Deligne-Mumford, the curve stable ordinaire $C_\eta$ degenere in such a way unique en a curve stable on $S$ (after ramification of $K$). It s'agit (after new ramification si requis) a faire degenerer $\mathcal{L}_\eta$ and the accouplement with values in $\omega_{C/S}$. En dehors of the points singuliers, one can prolonger $\mathcal{L}_\eta$ en a sheaf inversible $\mathcal{L}$. Corrigeant $\mathcal{L}$ by a diviseur concentre on the fiber speciale and, one can after passage a covering double of $S$ supposer that $\mathcal{L}^{\otimes 2} \to \omega$ en dehors of the points singuliers. After ramification, the direct image of this $\mathcal{L}$ conviendra.

To the neighborhood étale of each point singular, the racines carrees of $\omega_X$ on the complement of the point singular are classees by $H^1(\mu_2) = H^1(\mathbb{Z}/2)$: choisir a section $s$, and, in $\mathcal{L}$, on $(X-s)$, the $\sqrt{s}$ forment a covering double. The variety of the vanishing cycles a $H^1(\mathbb{Z}/2) = \mathbb{Z}/2$: after ramification on the base $S$ and localisation étale to the point singular, there is therefore deux cas possibles. They are couverts by the cas I and II of 2.3.

\vspace{1cm}

\begin{flushright}
Bien a toi,

\vspace{0.5cm}

P. Deligne
\end{flushright}

\end{document}
